\documentclass[a4paper,12pt]{article}

% ========================================================================
% Pacotes
% ========================================================================
\usepackage[margin=2.5cm]{geometry}
\usepackage[utf8]{inputenc}
\usepackage[T1]{fontenc}
\usepackage[brazil]{babel}
\usepackage{tikz}
\usepackage{amsmath}
\usepackage{amssymb}
\usepackage{graphicx}
\usepackage{float}
\usepackage{listings}
\usepackage{xcolor}
\usepackage{caption}
\usepackage{fancyhdr}

% TikZ
\usetikzlibrary{
    arrows.meta, positioning, calc, shapes.geometric,
    fit, decorations.markings, patterns
}

% ========================================================================
% Estilos TikZ compartilhados
% ========================================================================
\tikzset{
    block/.style = {
        rectangle, draw, very thick,
        minimum width=3.5cm, minimum height=1.8cm,
        font=\sffamily\small, align=center, text width=3cm
    },
    wide block/.style = {
        block, minimum width=5cm, text width=4.5cm
    },
    small block/.style = {
        rectangle, draw, thick,
        minimum width=2cm, minimum height=1cm,
        font=\sffamily\footnotesize, align=center
    },
    operator/.style = {
        circle, draw, very thick,
        minimum size=9mm, font=\sffamily\large,
        inner sep=0pt, fill=white
    },
    dot/.style = {
        circle, fill=black, minimum size=3pt, inner sep=0pt
    },
    port/.style = {font=\sffamily\bfseries\small},
    signal/.style = {font=\sffamily\footnotesize, fill=white, inner sep=1pt},
    bus/.style = {font=\sffamily\scriptsize},
    entity border/.style = {
        rectangle, draw, dashed, thick,
        rounded corners=3pt, inner sep=8pt
    },
    register/.style = {
        rectangle, draw, thick,
        minimum width=1.0cm, minimum height=0.8cm,
        font=\sffamily\footnotesize, align=center
    },
    data arrow/.style = {-Stealth, thick},
    ctrl arrow/.style = {-Stealth, thin, dashed},
    const/.style = {font=\sffamily\itshape\footnotesize},
    annotation/.style = {font=\sffamily\small, align=left},
    >=Stealth, thick, font=\sffamily
}

% ========================================================================
% Estilo VHDL
% ========================================================================
\definecolor{vhdlkeyword}{RGB}{0, 0, 160}
\definecolor{vhdlcomment}{RGB}{0, 128, 0}
\definecolor{vhdlbg}{RGB}{248, 248, 248}
\definecolor{vhdlnumber}{RGB}{120, 120, 120}

\lstdefinestyle{vhdlstyle}{
    language=VHDL,
    backgroundcolor=\color{vhdlbg},
    basicstyle=\ttfamily\scriptsize,
    keywordstyle=\color{vhdlkeyword}\bfseries,
    commentstyle=\color{vhdlcomment}\itshape,
    numberstyle=\tiny\color{vhdlnumber},
    numbers=left,
    numbersep=5pt,
    frame=single,
    rulecolor=\color{black!30},
    breaklines=true,
    tabsize=4,
    showstringspaces=false,
    captionpos=b,
    xleftmargin=2em,
    framexleftmargin=1.5em,
    morekeywords={data_t, data_long_t, data_array_t, signed, std_logic,
                  std_logic_vector, to_signed, shift_right, resize}
}

\lstdefinestyle{logstyle}{
    backgroundcolor=\color{black!5},
    basicstyle=\ttfamily\scriptsize,
    frame=single,
    rulecolor=\color{black!30},
    breaklines=true,
    captionpos=b,
    xleftmargin=1em,
    framexleftmargin=0.5em,
}

% ========================================================================
% Cabeçalho
% ========================================================================
\pagestyle{fancy}
\fancyhf{}
\fancyhead[L]{\sffamily\small Rede Neural DPD em VHDL}
\fancyhead[R]{\sffamily\small Atividade 12}
\fancyfoot[C]{\thepage}

% ========================================================================
% Título
% ========================================================================
\title{Relat\'orio -- Atividade 12 \\
       \large Implementa\c{c}\~ao das LUTs Trigonom\'etricas e Registrador de Deslocamento}
\author{Jo\~ao Pedro Hinning}
\date{\today}

% ========================================================================
% Documento
% ========================================================================
\begin{document}

\maketitle
\thispagestyle{fancy}

% =======================================================================
\section{Vis\~ao Geral}
% =======================================================================

Esta atividade documenta a implementa\c{c}\~ao e valida\c{c}\~ao de dois blocos fundamentais
do Est\'agio~1 da rede neural DPD: o \texttt{LUT\_Trig} e o \texttt{Shift\_Register}.
Ambos operam sobre dados em ponto fixo Q16.16 (\texttt{data\_t} = \texttt{signed(31 downto 0)}).

No pipeline do Est\'agio~1 (\texttt{Top\_NeuralNet\_S1}), o \texttt{Input\_Processing\_Block}
entrega o envelope $A(n)$ e a fase normalizada $\Delta\theta(n)$ ap\'os 40~ciclos de lat\^encia
(CORDIC). Os dois blocos desta atividade recebem esses sinais:

\begin{itemize}
    \item \textbf{\texttt{LUT\_Trig}:} Recebe $\Delta\theta(n)$ (\texttt{theta\_norm\_out}) e
          calcula $\cos(\Delta\theta)$ ou $\sin(\Delta\theta)$ em 4~ciclos de pipeline, usando
          uma tabela ROM de 1024~entradas com interpola\c{c}\~ao linear.
    \item \textbf{\texttt{Shift\_Register}:} Mant\'em uma janela deslizante de $M+1 = 6$
          amostras para cada canal (envelope, cosseno, seno), expondo simultaneamente todos os
          \textit{taps} para a camada convolucional subsequente.
\end{itemize}

O bloco de atraso de alinhamento presente no \texttt{Top\_NeuralNet\_S1} sincroniza o envelope
e a fase com as sa\'idas das LUTs antes de escrever nos registradores de deslocamento, garantindo
coer\^encia temporal entre os tr\^es canais.


% =======================================================================
\section{\texttt{LUT\_Trig} -- LUT Trigonom\'etrica com Pipeline}
% =======================================================================

\subsection{Arquitetura Interna}

O \texttt{LUT\_Trig} implementa o c\'alculo de $\cos(x)$ ou $\sin(x)$ para $x \in [-\pi, +\pi]$
em quatro est\'agios de pipeline registrados, com interpola\c{c}\~ao linear entre entradas
consecutivas da tabela. A Figura~\ref{fig:lut_pipeline} ilustra os est\'agios.

\begin{figure}[H]
\centering
\begin{tikzpicture}[node distance=0.6cm and 1.6cm, >=Stealth, thick, font=\sffamily]

\tikzset{
    stage/.style={
        rectangle, draw, very thick, fill=blue!8,
        minimum width=3.6cm, minimum height=2.6cm,
        font=\sffamily\small, align=center, text width=3.2cm
    },
    stagenum/.style={font=\sffamily\bfseries\normalsize},
    io/.style={font=\sffamily\bfseries\small},
    clkedge/.style={font=\sffamily\scriptsize, fill=yellow!30, draw=black, thick,
                    rectangle, minimum width=0.5cm, minimum height=0.4cm}
}

% Stage boxes
\node[stage] (S1) at (0,0) {
    \textbf{Est\'agio 1}\\[3pt]
    {\footnotesize C\'alculo de endere\c{c}o}\\[2pt]
    {\scriptsize $\text{addr} = x_{in} + \text{OFFSET}$}\\
    {\scriptsize $\text{idx} = \text{addr}[18\!:\!9]$}\\
    {\scriptsize $\text{addr\_norm\_d1} \leftarrow \text{addr}$}
};

\node[stage] (S2) at ($(S1.east)+(2.2,0)$) {
    \textbf{Est\'agio 2}\\[3pt]
    {\footnotesize Leitura da ROM}\\[2pt]
    {\scriptsize $y_0 \leftarrow \text{LUT}[\text{idx}]$}\\
    {\scriptsize $\text{slope} \leftarrow \text{SLOPES}[\text{idx}]$}\\
    {\scriptsize $\delta_x = \text{addr\_norm\_d1}$}\\
    {\scriptsize $\quad \mathbin{\&} \texttt{0x1FF}$}
};

\node[stage] (S3) at ($(S2.east)+(2.2,0)$) {
    \textbf{Est\'agio 3}\\[3pt]
    {\footnotesize Multiplica\c{c}\~ao}\\[2pt]
    {\scriptsize $\text{prod} = \delta_x \times \text{slope}$}\\[1pt]
    {\scriptsize (64~bits)}\\[2pt]
    {\scriptsize $y_{0\_d1} \leftarrow y_0$}\\
    {\scriptsize (atraso alinha\-mento)}
};

\node[stage] (S4) at ($(S3.east)+(2.2,0)$) {
    \textbf{Est\'agio 4}\\[3pt]
    {\footnotesize Soma Final}\\[2pt]
    {\scriptsize $\delta_y = \text{prod}[47\!:\!16]$}\\[1pt]
    {\scriptsize $y_{out} = y_{0\_d1} + \delta_y$}
};

% Input arrow
\node[io] (xin) at ($(S1.west)+(-2.0,0)$) {$x_{in}$ (\texttt{data\_t})};
\draw[data arrow] (xin) -- (S1.west);

% Arrows between stages
\draw[data arrow] (S1.east) -- (S2.west);
\draw[data arrow] (S2.east) -- (S3.west);
\draw[data arrow] (S3.east) -- (S4.west);

% Output arrow
\node[io] (yout) at ($(S4.east)+(2.0,0)$) {$y_{out}$ (\texttt{data\_t})};
\draw[data arrow] (S4.east) -- (yout);

% mode_sin annotation
\node[const, above=1.0cm of S2] (modeann) {\texttt{mode\_sin}: seleciona \texttt{LUT\_SIN\_Y/SLOPES} ou \texttt{LUT\_COS\_Y/SLOPES}};
\draw[ctrl arrow] (modeann) -- (S2.north);

% Clock ticks at bottom
\foreach \s/\n in {S1/1, S2/2, S3/3, S4/4} {
    \node[clkedge, below=0.5cm of \s] {CLK\,$\uparrow$\,\n};
}

\end{tikzpicture}
\caption{Pipeline de 4~est\'agios do \texttt{LUT\_Trig}. Cada est\'agio \'e registrado numa borda de subida do rel\'ogio.}
\label{fig:lut_pipeline}
\end{figure}

\subsection{Descri\c{c}\~ao dos Est\'agios}

\begin{enumerate}
    \item \textbf{Est\'agio~1 -- C\'alculo de endere\c{c}o:} A entrada $x_{in}$ \'e deslocada
          pelo \texttt{OFFSET\_VAL} $= 262144$ (que corresponde a $4{,}0$ em Q16.16, mapeando
          o intervalo $[-\pi,+\pi]$ para $[4{-}\pi,\, 4{+}\pi]$, ou seja, uma faixa sempre
          positiva). Os 10~bits $[18\!:\!9]$ do resultado formam o \'indice de 0 a 1023 que
          aponta para a entrada da tabela (\texttt{index}). O valor completo \'e preservado em
          \texttt{addr\_norm\_d1} para uso no Est\'agio~2.

    \item \textbf{Est\'agio~2 -- Leitura da ROM:} Com base em \texttt{mode\_sin}, s\~ao lidos
          \texttt{y0} e \texttt{slope} das constantes de pacote
          (\texttt{LUT\_SIN\_Y}/\texttt{LUT\_SIN\_SLOPES} ou
          \texttt{LUT\_COS\_Y}/\texttt{LUT\_COS\_SLOPES}). Simultaneamente, $\delta_x$ \'e
          obtido mascarando os 9~bits inferiores de \texttt{addr\_norm\_d1}
          (\texttt{\& 0x1FF}), representando a fra\c{c}\~ao dentro do intervalo da entrada da tabela.

    \item \textbf{Est\'agio~3 -- Multiplica\c{c}\~ao e alinhamento:} O produto
          $\text{prod\_long} = \delta_x \times \text{slope}$ \'e calculado em 64~bits para
          preservar precis\~ao. Como a leitura ROM e a multiplica\c{c}\~ao est\~ao em est\'agios
          separados, \texttt{y0} \'e atrasado um ciclo em \texttt{y0\_d1} para permanecer
          alinhado com o resultado da multiplica\c{c}\~ao.

    \item \textbf{Est\'agio~4 -- Soma final:} A corre\c{c}\~ao linear $\delta_y =
          \text{prod\_long}[47\!:\!16]$ (extrai\c{c}\~ao dos bits de precis\~ao Q16.16 do
          produto de 64~bits) \'e somada a \texttt{y0\_d1} para produzir \texttt{y\_out}.
\end{enumerate}

\subsection{Tabela ROM e Interpola\c{c}\~ao}

A LUT possui \textbf{1024~entradas} cobrindo o intervalo $[-\pi,\,+\pi]$. O deslocamento de
\texttt{OFFSET\_VAL} $= 262144$ ($= 4{,}0$ em Q16.16) garante que o argumento somado seja
sempre positivo, permitindo a indexa\c{c}\~ao direta via bits inteiros. A interpola\c{c}\~ao
linear reduz o tamanho necess\'ario de tabela: em vez de armazenar valores em alta resolu\c{c}\~ao
de fase, basta armazenar $y_0$ e a inclina\c{c}\~ao local (\texttt{slope}) para cada intervalo,
e calcular $y = y_0 + \delta_x \times \text{slope}$ com precis\~ao equivalente a uma tabela muito
maior. O sinal \texttt{mode\_sin} seleciona qual par de tabelas (seno ou cosseno) \'e utilizado.

% =======================================================================
\section{\texttt{Shift\_Register} -- Registrador de Deslocamento}
% =======================================================================

\subsection{Arquitetura}

O \texttt{Shift\_Register} implementa uma janela deslizante (\textit{sliding window}) de
$\text{DEPTH}+1$ posi\c{c}\~oes, parametriz\'avel via gen\'erico \texttt{DEPTH}. No projeto,
\texttt{DEPTH} $= M = 5$ resulta em um buffer de 6~taps. A Figura~\ref{fig:shift_reg} mostra
o mecanismo de deslocamento para \texttt{DEPTH} $= 5$.

\begin{figure}[H]
\centering
\begin{tikzpicture}[>=Stealth, thick, font=\sffamily]

\tikzset{
    cell/.style={rectangle, draw, very thick, minimum width=1.5cm, minimum height=1.0cm,
                 font=\sffamily\small, align=center, fill=green!8},
    io/.style={font=\sffamily\bfseries\small},
    idx/.style={font=\sffamily\scriptsize\itshape}
}

% Cells
\foreach \i/\lbl in {0/regs(0), 1/regs(1), 2/regs(2), 3/regs(3), 4/regs(4), 5/regs(5)} {
    \node[cell] (R\i) at (\i*1.8, 0) {\lbl};
    \node[idx] at (\i*1.8, -0.8) {tap \i};
}

% data_in arrow
\node[io] (din) at (-2.2, 0) {\texttt{data\_in}};
\draw[data arrow] (din) -- (R0.west) node[signal,midway,above]{\scriptsize \texttt{enable}};

% data_out arrows
\foreach \i in {0,1,2,3,4,5} {
    \draw[-Stealth, thick, blue!70] (R\i.south) -- ++(0,-0.9)
        node[below, font=\sffamily\scriptsize] {\texttt{data\_out(\i)}};
}

% shift arrows between cells
\foreach \i/\j in {0/1, 1/2, 2/3, 3/4, 4/5} {
    \draw[data arrow, dashed, red!70] (R\i.east) -- (R\j.west)
        node[signal,midway,above]{\scriptsize $\leftarrow$};
}

% enable annotation
\node[font=\sffamily\small, above=1.2cm of R2, align=center]
    {Quando \texttt{enable = '1'}:\\[2pt]
     \texttt{regs(0) <= data\_in}\\
     \texttt{regs(i) <= regs(i-1)}, $i = 1\ldots\text{DEPTH}$};

% hold annotation
\node[font=\sffamily\small, below=2.2cm of R4, align=center, text=red!70]
    {Quando \texttt{enable = '0'}: buffer congela (Hold)};

\end{tikzpicture}
\caption{Janela deslizante do \texttt{Shift\_Register} com \texttt{DEPTH} $= 5$ ($M = 5$, buffer de 6~taps). As setas vermelhas tracejadas indicam o deslocamento a cada pulso \texttt{enable}.}
\label{fig:shift_reg}
\end{figure}

\subsection{Gen\'erico \texttt{DEPTH}}

O par\^ametro gen\'erico \texttt{DEPTH} define o tamanho da mem\'oria (n\'umero de amostras
antigas armazenadas). O buffer ter\'a \texttt{DEPTH}$+1$~posi\c{c}\~oes no total, expondo
todos os \textit{taps} simultaneamente na porta \texttt{data\_out}. Isso permite que a camada
convolucional acesse toda a janela de hist\'orico sem lat\^encia adicional de leitura. No
\texttt{Top\_NeuralNet\_S1}, os tr\^es registradores s\~ao instanciados com
\texttt{Generic Map(M)} onde $M = 5$, produzindo buffers de 6~elementos.

Quando \texttt{enable = '0'}, o processo n\~ao executa e todos os registradores permanecem
congelados, implementando a fun\c{c}\~ao de \textit{hold} necess\'aria para respeitar o controle
de fluxo do pipeline.


% =======================================================================
\section{C\'odigo VHDL}
% =======================================================================

\subsection{\texttt{LUT\_Trig}}

\begin{lstlisting}[style=vhdlstyle,
    caption={C\'odigo VHDL do \texttt{LUT\_Trig}.},
    label={lst:lut_trig}]
library IEEE;
use IEEE.STD_LOGIC_1164.ALL;
use IEEE.NUMERIC_STD.ALL;
use work.pkg_neural_types.ALL;
use work.pkg_model_constants.ALL;

entity LUT_Trig is
    Port (
        clk    : in  std_logic;
        rst    : in  std_logic;
        x_in   : in  data_t;
        mode_sin : in std_logic;
        y_out  : out data_t
    );
end LUT_Trig;

architecture Behavioral of LUT_Trig is

    signal index : integer range 0 to 1023 := 0;
    constant OFFSET_VAL : signed(31 downto 0) := to_signed(262144, 32);

    -- Sinais de Estagio 1
    signal addr_norm_d1 : signed(31 downto 0) := (others => '0');

    -- Sinais de Estagio 2
    signal y0, slope : data_t := (others => '0');
    signal delta_x   : data_t := (others => '0');

    -- Sinais de Estagio 3 (Pipeline para sincronia)
    signal y0_d1     : data_t := (others => '0'); -- Delay para y0
    signal prod_long : signed(63 downto 0) := (others => '0');

    -- Sinais de Estagio 4
    signal delta_y   : data_t := (others => '0');

begin

    process(clk)
        variable address_norm : signed(31 downto 0);
        variable idx_calc : integer;
    begin
        if rising_edge(clk) then
            if rst = '1' then
                y_out <= (others => '0');
                index <= 0;
                addr_norm_d1 <= (others => '0');
                y0 <= (others => '0');
                slope <= (others => '0');
                delta_x <= (others => '0');
                y0_d1 <= (others => '0');
                prod_long <= (others => '0');
            else
                -- === ESTAGIO 1: Calculo de Endereco ===
                address_norm := x_in + OFFSET_VAL;

                -- Slice e Cast Unsigned
                idx_calc := to_integer(unsigned(address_norm(18 downto 9)));

                if idx_calc < 0 then index <= 0;
                elsif idx_calc > 1023 then index <= 1023;
                else index <= idx_calc;
                end if;

                addr_norm_d1 <= address_norm;

                -- === ESTAGIO 2: Leitura ROM e Delta X ===
                if mode_sin = '1' then
                    y0 <= LUT_SIN_Y(index);
                    slope <= LUT_SIN_SLOPES(index);
                else
                    y0 <= LUT_COS_Y(index);
                    slope <= LUT_COS_SLOPES(index);
                end if;

                delta_x <= addr_norm_d1 and x"000001FF";

                -- === ESTAGIO 3: Multiplicacao e Delay de Alinhamento ===
                prod_long <= delta_x * slope;

                -- Atrasar y0 para alinhar com o resultado da mult
                y0_d1 <= y0;

                -- === ESTAGIO 4: Soma Final ===
                delta_y <= prod_long(47 downto 16);
                y_out <= y0_d1 + delta_y;

            end if;
        end if;
    end process;

end Behavioral;
\end{lstlisting}

\subsection{\texttt{Shift\_Register}}

\begin{lstlisting}[style=vhdlstyle,
    caption={C\'odigo VHDL do \texttt{Shift\_Register}.},
    label={lst:shift_reg}]
library IEEE;
use IEEE.STD_LOGIC_1164.ALL;
use IEEE.NUMERIC_STD.ALL;
use work.pkg_neural_types.ALL;

entity Shift_Register is
    Generic (
        DEPTH : integer := 4 -- Tamanho da memoria (M). O buffer tera M+1 posicoes.
    );
    Port (
        clk      : in  STD_LOGIC;
        rst      : in  STD_LOGIC;
        enable   : in  STD_LOGIC; -- So desloca quando uma nova amostra valida chega

        -- Entrada: A nova amostra x(n)
        data_in  : in  data_t;

        -- Saida: O vetor completo [x(n), x(n-1), ..., x(n-M)]
        -- Permite que a Convolucao acesse TODOS os taps simultaneamente
        data_out : out data_array_t(0 to DEPTH)
    );
end Shift_Register;

architecture Behavioral of Shift_Register is
    -- Sinal interno para armazenar os valores (os registradores fisicos)
    signal regs : data_array_t(0 to DEPTH) := (others => (others => '0'));
begin

    process(clk)
    begin
        if rising_edge(clk) then
            if rst = '1' then
                -- Reset sincrono: Zera tudo
                regs <= (others => (others => '0'));
            elsif enable = '1' then
                -- MECANICA DE DESLOCAMENTO (SHIFT)
                -- 1. Insere a nova amostra na cabeca (indice 0)
                regs(0) <= data_in;

                -- 2. Move os antigos para a direita
                for i in 1 to DEPTH loop
                    regs(i) <= regs(i-1);
                end loop;
            end if;
        end if;
    end process;

    -- Conecta os registradores internos a porta de saida
    data_out <= regs;

end Behavioral;
\end{lstlisting}


% =======================================================================
\section{Testbench e Resultados}
% =======================================================================

\subsection{Testbench do \texttt{LUT\_Trig}}

\begin{lstlisting}[style=vhdlstyle,
    caption={Testbench do \texttt{LUT\_Trig}.},
    label={lst:lut_trig_tb}]
library IEEE;
use IEEE.STD_LOGIC_1164.ALL;
use IEEE.NUMERIC_STD.ALL;
use work.pkg_neural_types.ALL;

entity LUT_Trig_tb is
end LUT_Trig_tb;

architecture Behavioral of LUT_Trig_tb is

    component LUT_Trig
    Port (
        clk    : in  std_logic;
        rst    : in  std_logic;
        x_in   : in  data_t;
        mode_sin : in std_logic;
        y_out  : out data_t
    );
    end component;

    signal clk : std_logic := '0';
    signal rst : std_logic := '0';
    signal x_in : data_t := (others => '0');
    signal mode_sin : std_logic := '0';
    signal y_out : data_t;

    constant clk_period : time := 10 ns;

    -- Constants Q16.16
    -- Pi approx 3.14159 * 65536 = 205887
    constant PI_VAL  : integer := 205887;
    constant HALF_PI : integer := 102943;
    constant VAL_ONE : integer := 65536;

begin

    uut: LUT_Trig PORT MAP (
        clk => clk,
        rst => rst,
        x_in => x_in,
        mode_sin => mode_sin,
        y_out => y_out
    );

    clk_process :process
    begin
        clk <= '0';
        wait for clk_period/2;
        clk <= '1';
        wait for clk_period/2;
    end process;

    stim_proc: process
    begin
        rst <= '1';
        wait for 100 ns;
        rst <= '0';
        wait for clk_period*2;

        report "--- STARTING LUT_TRIG TEST ---";

        -- === TEST COSINE (Mode = 0) ===
        mode_sin <= '0';

        -- Case 1: Cos(0) -> 1.0 (65536)
        x_in <= to_signed(0, 32);
        wait for clk_period * 5;
        report "Cos(0): " & integer'image(to_integer(y_out)) & " (Exp: 65536)";
        assert abs(to_integer(y_out) - VAL_ONE) < 50 severity warning;

        -- Case 2: Cos(Pi/2) -> 0.0
        x_in <= to_signed(HALF_PI, 32);
        wait for clk_period * 5;
        report "Cos(Pi/2): " & integer'image(to_integer(y_out)) & " (Exp: 0)";
        assert abs(to_integer(y_out)) < 100 severity warning;

        -- Case 3: Cos(Pi) -> -1.0 (-65536)
        x_in <= to_signed(PI_VAL, 32);
        wait for clk_period * 5;
        report "Cos(Pi): " & integer'image(to_integer(y_out)) & " (Exp: -65536)";
        assert abs(to_integer(y_out) - (-VAL_ONE)) < 50 severity warning;

        -- === TEST SINE (Mode = 1) ===
        mode_sin <= '1';

        -- Case 4: Sin(0) -> 0.0
        x_in <= to_signed(0, 32);
        wait for clk_period * 5;
        report "Sin(0): " & integer'image(to_integer(y_out)) & " (Exp: 0)";
        assert abs(to_integer(y_out)) < 50 severity warning;

        -- Case 5: Sin(Pi/2) -> 1.0
        x_in <= to_signed(HALF_PI, 32);
        wait for clk_period * 5;
        report "Sin(Pi/2): " & integer'image(to_integer(y_out)) & " (Exp: 65536)";
        assert abs(to_integer(y_out) - VAL_ONE) < 50 severity warning;

        -- Case 6: Sin(-Pi/2) -> -1.0
        x_in <= to_signed(-HALF_PI, 32);
        wait for clk_period * 5;
        report "Sin(-Pi/2): " & integer'image(to_integer(y_out)) & " (Exp: -65536)";
        assert abs(to_integer(y_out) - (-VAL_ONE)) < 50 severity warning;

        report "--- END OF TEST ---";
        wait;
    end process;

end Behavioral;
\end{lstlisting}

\subsection{Resultados do ISim -- \texttt{LUT\_Trig}}

\begin{figure}[H]
\centering
\fbox{\parbox{0.88\textwidth}{\centering\vspace{2.0cm}
    {\sffamily\color{black!45} [Captura de tela do ISim -- LUT\_Trig]}\\[4pt]
    {\sffamily\small\color{black!45} Salvar como \texttt{isim\_lut\_trig.png} e descomentar \textbackslash includegraphics}
\vspace{2.0cm}}}
\caption{Formas de onda da simula\c{c}\~ao do \texttt{LUT\_Trig} no ISim.}
\label{fig:wave_lut}
\end{figure}

\begin{lstlisting}[style=logstyle,
    caption={Log de sa\'ida do ISim -- \texttt{LUT\_Trig}.},
    label={lst:log_lut}]
Finished circuit initialization process.
at 100 ns:  Note: --- STARTING LUT_TRIG TEST ---
at 170 ns:  Note: Cos(0): 65536 (Exp: 65536)
at 220 ns:  Note: Cos(Pi/2): 3 (Exp: 0)
at 270 ns:  Note: Cos(Pi): -65533 (Exp: -65536)
at 330 ns:  Note: Sin(0): 1 (Exp: 0)
at 380 ns:  Note: Sin(Pi/2): 65536 (Exp: 65536)
at 430 ns:  Note: Sin(-Pi/2): -65536 (Exp: -65536)
at 430 ns:  Note: --- END OF TEST ---
\end{lstlisting}

\begin{table}[H]
\centering
\caption{Resultados quantitativos do \texttt{LUT\_Trig}.}
\label{tab:lut_results}
\small
\begin{tabular}{llccc}
\hline
\textbf{Caso} & \textbf{Fun\c{c}\~ao} & \textbf{Esperado} & \textbf{Obtido} & \textbf{Erro (LSB)} \\
\hline
$\cos(0)$          & Cosseno & 65536  & 65536  & 0   \\
$\cos(\pi/2)$      & Cosseno & 0      & 3      & +3  \\
$\cos(\pi)$        & Cosseno & -65536 & -65533 & +3  \\
$\sin(0)$          & Seno    & 0      & 1      & +1  \\
$\sin(\pi/2)$      & Seno    & 65536  & 65536  & 0   \\
$\sin(-\pi/2)$     & Seno    & -65536 & -65536 & 0   \\
\hline
\end{tabular}
\end{table}

Todos os resultados est\~ao dentro da toler\^ancia de $\pm 50$~LSB definida no testbench.
Os erros residuais de 1--3~LSB s\~ao inerentes \`a quantiza\c{c}\~ao Q16.16 e \`a
interpola\c{c}\~ao linear com 1024~entradas.

\subsection{Testbench do \texttt{Shift\_Register}}

\begin{lstlisting}[style=vhdlstyle,
    caption={Testbench do \texttt{Shift\_Register} (DEPTH = 4, buffer de 5 posi\c{c}\~oes).},
    label={lst:sr_tb}]
library IEEE;
use IEEE.STD_LOGIC_1164.ALL;
use IEEE.NUMERIC_STD.ALL;
use work.pkg_neural_types.ALL;

entity Shift_Register_tb is
end Shift_Register_tb;

architecture Behavioral of Shift_Register_tb is

    component Shift_Register
        Generic ( DEPTH : integer );
        Port (
            clk      : in  STD_LOGIC;
            rst      : in  STD_LOGIC;
            enable   : in  STD_LOGIC;
            data_in  : in  data_t;
            data_out : out data_array_t(0 to DEPTH)
        );
    end component;

    signal clk : std_logic := '0';
    signal rst : std_logic := '0';
    signal enable : std_logic := '0';
    signal data_in : data_t := (others => '0');

    constant TEST_DEPTH : integer := 4;
    signal data_out : data_array_t(0 to TEST_DEPTH);

    constant clk_period : time := 10 ns;

begin

    uut: Shift_Register
    Generic Map ( DEPTH => TEST_DEPTH )
    Port Map (
        clk => clk, rst => rst,
        enable => enable,
        data_in => data_in,
        data_out => data_out
    );

    clk_process :process
    begin
        clk <= '0'; wait for clk_period/2;
        clk <= '1'; wait for clk_period/2;
    end process;

    stim_proc: process
    begin
        -- TESTE 1: RESET
        report "Iniciando Teste: Reset";
        rst <= '1'; wait for 20 ns; rst <= '0'; wait for 10 ns;
        assert to_integer(data_out(0)) = 0 report "Falha no Reset" severity error;

        -- TESTE 2: SHIFT (Encher o buffer)
        report "Iniciando Teste: Shift Enable";
        enable <= '1';
        data_in <= to_signed(10, 32); wait for clk_period;  -- [10,0,0,0,0]
        data_in <= to_signed(20, 32); wait for clk_period;  -- [20,10,0,0,0]
        data_in <= to_signed(30, 32); wait for clk_period;  -- [30,20,10,0,0]
        data_in <= to_signed(40, 32); wait for clk_period;  -- [40,30,20,10,0]
        data_in <= to_signed(50, 32); wait for clk_period;  -- [50,40,30,20,10]

        assert to_integer(data_out(0)) = 50 report "Erro pos 0" severity error;
        assert to_integer(data_out(1)) = 40 report "Erro pos 1" severity error;
        assert to_integer(data_out(TEST_DEPTH)) = 10 report "Erro pos final" severity error;

        -- TESTE 3: HOLD (Enable = 0)
        report "Iniciando Teste: Hold (Enable=0)";
        enable <= '0';
        data_in <= to_signed(999, 32); -- Valor lixo, nao deve entrar
        wait for clk_period * 3;
        assert to_integer(data_out(0)) = 50 report "Erro Hold: Dado mudou!" severity error;

        -- TESTE 4: RESUME
        report "Iniciando Teste: Resume Shift";
        enable <= '1';
        data_in <= to_signed(60, 32);
        wait for clk_period;  -- [60,50,40,30,20]

        assert to_integer(data_out(0)) = 60 report "Erro Resume" severity error;
        assert to_integer(data_out(1)) = 50 report "Erro Deslocamento Resume" severity error;

        report "--- FIM DO TESTE: SUCESSO ---";
        wait;
    end process;

end Behavioral;
\end{lstlisting}

\subsection{Resultados do ISim -- \texttt{Shift\_Register}}

\begin{figure}[H]
\centering
\fbox{\parbox{0.88\textwidth}{\centering\vspace{2.0cm}
    {\sffamily\color{black!45} [Captura de tela do ISim -- Shift\_Register]}\\[4pt]
    {\sffamily\small\color{black!45} Salvar como \texttt{isim\_shift\_register.png} e descomentar \textbackslash includegraphics}
\vspace{2.0cm}}}
\caption{Formas de onda da simula\c{c}\~ao do \texttt{Shift\_Register} no ISim.}
\label{fig:wave_sr}
\end{figure}

\begin{lstlisting}[style=logstyle,
    caption={Log de sa\'ida do ISim -- \texttt{Shift\_Register}.},
    label={lst:log_sr}]
Finished circuit initialization process.
at  10 ns:  Note: Iniciando Teste: Reset
at  40 ns:  Note: Iniciando Teste: Shift Enable
at 100 ns:  -- data_out = [10, 0,  0,  0,  0]
at 110 ns:  -- data_out = [20, 10, 0,  0,  0]
at 120 ns:  -- data_out = [30, 20, 10, 0,  0]
at 130 ns:  -- data_out = [40, 30, 20, 10, 0]
at 140 ns:  -- data_out = [50, 40, 30, 20, 10]  (Buffer Cheio)
at 140 ns:  Note: Iniciando Teste: Hold (Enable=0)
at 170 ns:  -- data_out = [50, 40, 30, 20, 10]  (INALTERADO)
at 170 ns:  Note: Iniciando Teste: Resume Shift
at 180 ns:  -- data_out = [60, 50, 40, 30, 20]
at 180 ns:  Note: --- FIM DO TESTE: SUCESSO ---
\end{lstlisting}

\begin{table}[H]
\centering
\caption{Resumo dos testes do \texttt{Shift\_Register}.}
\label{tab:sr_results}
\small
\begin{tabular}{llcc}
\hline
\textbf{Teste} & \textbf{Descri\c{c}\~ao} & \textbf{Resultado} & \textbf{Status} \\
\hline
1 -- Reset       & Todos os taps zerados ap\'os reset s\'incrono        & \texttt{data\_out} $= [0,0,0,0,0]$ & PASS \\
2 -- Shift       & Buffer cheio com 5 amostras sequenciais (10\ldots50) & \texttt{[50,40,30,20,10]}           & PASS \\
3 -- Hold        & Enable=0 por 3 ciclos: buffer n\~ao altera           & \texttt{[50,40,30,20,10]}           & PASS \\
4 -- Resume      & Retomada com dado 60: desloca corretamente           & \texttt{[60,50,40,30,20]}           & PASS \\
\hline
\end{tabular}
\end{table}


% =======================================================================
\section{Pr\'oximos Passos}
% =======================================================================

\begin{enumerate}
    \item Integrar o \texttt{LUT\_Trig} e o \texttt{Shift\_Register} com o
          \texttt{Input\_Processing\_Block} no \texttt{Top\_NeuralNet\_S1}, incluindo o
          bloco de atraso de alinhamento de 4~ciclos.
    \item Validar o Est\'agio~1 completo com vetores de teste extra\'idos de refer\^encia
          Python (\texttt{vectors\_s1\_input.txt}).
    \item Implementar a camada convolucional (\texttt{Conv\_Layer\_Serial}) como pr\'oximo
          est\'agio do pipeline.
\end{enumerate}

\end{document}
